\section{11. Conclusion}

\subsection{11.1 Summary}

We presented \textbf{ClusterFormer}, a differentiable K-means clustering 
attention mechanism that achieves $\mathcal{O}(KLd)$ complexity while 
preserving global context. The method combines centroid-to-key 
cross-attention with local sliding-window blending, controlled by learnable 
temperature and blend scalars.

\subsection{11.2 Results Recap}

\begin{table}[H]
\centering
\begin{tabular}{lll}
\toprule
\textbf{Model} & \textbf{Acc Δ} & \textbf{FLOPs reduction ($L=256$)} \\
\midrule
DistilBERT & -0.80\% & ~2.7× \\
RoBERTa & -1.06\% & ~2.7× \\
GPT-2 & -0.38\% & ~2.7× \\
\bottomrule
\end{tabular}
\end{table}

\subsection{11.3 Key Takeaways}

1. \textbf{Linear attention without random features} is possible via clustering.
2. \textbf{Global + local blending} is essential.
3. \textbf{Learnable hyperparameters} (τ, α) improve results marginally.
4. \textbf{Not a compression method} — parameter count increases.

\subsection{11.4 Future Work}

\subsubsection{Short-Term}
\begin{itemize}
\item Fused Triton kernel for K-means
\end{itemize}
\begin{itemize}
\item Learned K per layer
\end{itemize}
\begin{itemize}
\item Long-context evaluation ($L \geq 4096$)
\end{itemize}

\subsubsection{Long-Term}
\begin{itemize}
\item Causal generation support
\end{itemize}
\begin{itemize}
\item Theory: bound on approximation error
\end{itemize}
\begin{itemize}
\item Combination with quantization
\end{itemize}

\subsection{11.5 Closing Remarks}

ClusterFormer is a practical linear attention method for long sequences. 
We release code and checkpoints under MIT.

\subsection{11.6 Acknowledgments}

{{ACKNOWLEDGMENTS}}

\subsection{11.7 Author Contributions}

\begin{itemize}
\item \textbf{{{AUTHOR_1}}} designed the method, implemented the code, ran experiments.
\end{itemize}
\begin{itemize}
\item \textbf{{{AUTHOR_2}}} provided feedback and edited the manuscript.
\end{itemize}

\subsection{11.8 Conflict of Interest}

None.